\documentclass[11pt, a4paper]{article}

% Gemeinsame Style-Definitionen (TEXINPUTS muss templates/ enthalten — siehe scripts/build_tex.py)
% bfw-style.tex — gemeinsame Style-Definitionen für alle Handouts/Aufgabenhefte
% Voraussetzung: Kompilation mit xelatex (wegen fontspec)

\usepackage[ngerman]{babel}
\usepackage{geometry}
\geometry{a4paper, top=2.2cm, bottom=2.2cm, left=2.3cm, right=2.3cm}

\usepackage{fontspec}
% Fallback-Kette: Source Sans Pro -> Calibri -> Segoe UI -> Latin Modern Sans
% (Latin Modern ist immer mit MiKTeX vorhanden)
\IfFontExistsTF{Source Sans Pro}{%
  \setmainfont{Source Sans Pro}[Ligatures=TeX]%
}{\IfFontExistsTF{Calibri}{%
  \setmainfont{Calibri}[Ligatures=TeX]%
}{\IfFontExistsTF{Segoe UI}{%
  \setmainfont{Segoe UI}[Ligatures=TeX]%
}{%
  \setmainfont{Latin Modern Sans}[Ligatures=TeX]%
}}}
\IfFontExistsTF{Source Code Pro}{%
  \setmonofont{Source Code Pro}[Scale=0.92]%
}{\IfFontExistsTF{Consolas}{%
  \setmonofont{Consolas}[Scale=0.92]%
}{%
  \setmonofont{Latin Modern Mono}[Scale=0.92]%
}}

\usepackage{xcolor}
% Farbpalette: hell, freundlich, modern
\definecolor{bfwPrimary}{HTML}{0EA5A4}    % Petrol/Türkis - Primär
\definecolor{bfwPrimaryDark}{HTML}{0F766E} % dunkler Petrol für Headings
\definecolor{bfwAccent}{HTML}{F59E0B}     % Amber - Akzent
\definecolor{bfwSuccess}{HTML}{10B981}    % Grün - Erfolg / Lösung
\definecolor{bfwWarn}{HTML}{F97316}       % Orange - Achtung
\definecolor{bfwInk}{HTML}{0F172A}        % fast Schwarz
\definecolor{bfwInkSoft}{HTML}{475569}    % gedämpftes Grau
\definecolor{bfwBgSoft}{HTML}{F8FAFC}     % off-white
\definecolor{bfwBgTeal}{HTML}{ECFEFF}     % helles Türkis
\definecolor{bfwBgAmber}{HTML}{FEF3C7}    % helles Gelb
\definecolor{bfwBgRose}{HTML}{FFE4E6}     % helles Rosé für Eselsbrücken

\usepackage[colorlinks=true, linkcolor=bfwPrimaryDark, urlcolor=bfwPrimaryDark]{hyperref}
\usepackage{enumitem}
\setlist[itemize]{leftmargin=*, itemsep=2pt, topsep=2pt}
\setlist[enumerate]{leftmargin=*, itemsep=2pt, topsep=2pt}

\usepackage[most]{tcolorbox}
\usepackage{booktabs}
\usepackage{tikz}
\usetikzlibrary{decorations.pathmorphing, positioning, shapes.geometric, arrows.meta}

% Merkkästchen — wichtige Definition / Kernpunkt
\newtcolorbox{merksatz}[1][]{%
  enhanced, breakable,
  colback=bfwBgTeal,
  colframe=bfwPrimary,
  boxrule=0.6pt,
  arc=4pt,
  left=10pt, right=10pt, top=8pt, bottom=8pt,
  fonttitle=\bfseries\color{bfwPrimaryDark},
  title={\faLightbulb\ Merksatz},
  #1
}

% Eselsbrücke — Mnemoniks
\newtcolorbox{eselsbruecke}[1][]{%
  enhanced, breakable,
  colback=bfwBgRose,
  colframe=bfwAccent,
  boxrule=0.6pt,
  arc=4pt,
  left=10pt, right=10pt, top=8pt, bottom=8pt,
  fonttitle=\bfseries\color{bfwWarn},
  title={Eselsbrücke},
  #1
}

% Beispiel-Box
\newtcolorbox{beispiel}[1][]{%
  enhanced, breakable,
  colback=bfwBgSoft,
  colframe=bfwInkSoft,
  boxrule=0.4pt,
  arc=3pt,
  left=10pt, right=10pt, top=8pt, bottom=8pt,
  fonttitle=\bfseries\color{bfwInk},
  title={Beispiel},
  #1
}

% Lösungs-Box (für Aufgabenhefte)
\newtcolorbox{loesung}[1][]{%
  enhanced, breakable,
  colback=bfwBgAmber!40,
  colframe=bfwSuccess,
  boxrule=0.6pt,
  arc=3pt,
  left=10pt, right=10pt, top=8pt, bottom=8pt,
  fonttitle=\bfseries\color{bfwSuccess!70!black},
  title={Lösung},
  #1
}

% Glossar-Eintrag
\newcommand{\glossareintrag}[2]{%
  \par\noindent\textbf{\color{bfwPrimaryDark}#1} \enspace #2\par\smallskip
}

% Section-Stil — etwas farbiger als Standard
\usepackage{titlesec}
\titleformat{\section}
  {\Large\bfseries\color{bfwPrimaryDark}}
  {\thesection}{0.6em}{}
\titleformat{\subsection}
  {\large\bfseries\color{bfwPrimary}}
  {\thesubsection}{0.5em}{}
\titleformat{\subsubsection}
  {\normalsize\bfseries\color{bfwInk}}
  {\thesubsubsection}{0.4em}{}

% TikZ-Stil "skizzenhaft" — etwas weicher als Rough.js, aber stabil
\tikzset{
  roughbox/.style = {
    draw=bfwPrimaryDark, thick, fill=bfwBgTeal,
    rounded corners=4pt, inner sep=6pt
  },
  roughaccent/.style = {
    draw=bfwAccent, thick, fill=bfwBgAmber,
    rounded corners=4pt, inner sep=6pt
  },
  roughline/.style = {
    draw=bfwInkSoft, thick, -{Stealth[length=2.5mm]}
  }
}

% Bullet-Symbole (bewusst kein FontAwesome — vermeidet Font-Installations-Probleme)
\newcommand{\faLightbulb}{$\bullet$}
\newcommand{\faBookmark}{$\bullet$}

% Header/Footer
\usepackage{fancyhdr}
\pagestyle{fancy}
\fancyhf{}
\renewcommand{\headrulewidth}{0pt}
\fancyfoot[L]{\small\color{bfwInkSoft}\thepage}
\fancyfoot[R]{\small\color{bfwInkSoft} BFW M-064 Beschaffung im Büromanagement}


\title{\vspace*{-1.5cm}\Huge\bfseries\color{bfwPrimaryDark}Tag~7: Sitzungen \& Besprechungen vorbereiten und durchführen\\[0.4em]
  \large\color{bfwInkSoft}\textnormal{Wie gute Planung aus Meetings echte Ergebnisse macht}}
\author{\textsf{Büroprozesse gestalten und Arbeitsvorgänge organisieren \textperiodcentered{} M-067}\\
\textsf{\small\copyright\ Can Siebert\,\textperiodcentered\,2026}}
\date{07.07.2026}

\begin{document}
\maketitle
\thispagestyle{empty}

\vspace{1em}
\begin{merksatz}[title={\bfseries Worum geht es heute?}]
Besprechungen beanspruchen einen erheblichen Teil der Arbeitszeit im Büro --- und nichts
ist ärgerlicher als ein Meeting, das nichts bringt. Ob eine Sitzung zufriedenstellend und
effektiv verläuft, entscheidet sich meist schon \emph{vor} dem ersten Wort: bei der
Vorbereitung. In diesem Handout lernen Sie, wie Sie Ziel und Teilnehmerkreis bestimmen,
eine präzise Tagesordnung und eine vollständige Einladung erstellen, den zeitlichen Ablauf
der Vorbereitung planen und eine Besprechung souverän durchführen --- inklusive
Rollenverteilung, Diskussionsregeln und moderner Video- und Hybridkonferenzen.
\end{merksatz}

\tableofcontents
\vspace{1em}
\hrule
\vspace{1em}

\section{Warum gute Vorbereitung über den Erfolg entscheidet}

Sollen Sitzungen und Besprechungen zufriedenstellend und effektiv verlaufen, ist eine gute
Planung \emph{unverzichtbar}. Das gilt auch für Teambesprechungen mit kleiner
Teilnehmerzahl. Bei mangelnder Vorbereitung droht gleich eine ganze Kette von Problemen:
Wichtige Inhalte werden nicht thematisiert, die Sitzung beginnt verspätet und benötigt mehr
Zeit als geplant, und Ergebnisse oder Beschlüsse werden möglicherweise nicht genau
festgehalten, nicht nachverfolgt und schließlich vergessen.

\begin{merksatz}
Eine schlecht vorbereitete Besprechung kostet doppelt: einmal die verschwendete Zeit aller
Anwesenden und ein zweites Mal, weil die offenen Punkte später erneut bearbeitet werden
müssen. Vorbereitung ist daher gelebte Wirtschaftlichkeit.
\end{merksatz}

\subsection{Rechtzeitig beginnen --- und Fristen beachten}

Eine gute Vorbereitung beginnt \emph{rechtzeitig}. Wie groß der Zeitabstand zwischen
Vorbereitungsbeginn und Sitzungstermin sein muss, hängt von der Art der Sitzung und den
konkreten Bedingungen vor Ort ab. Für eine wöchentlich stattfindende Teamsitzung reicht es
in der Regel aus, mit der Planung 2--3 Tage vorher zu beginnen. Handelt es sich dagegen um
eine \emph{offizielle} Sitzung wie die Gesellschafterversammlung einer GmbH, sind die
gesetzlichen Regeln einzuhalten und die Vorbereitung sollte entsprechend früher beginnen.

\begin{beispiel}
\textbf{GmbH-Gesetz § 51 (Form der Einberufung):} Die Berufung der Versammlung erfolgt durch
Einladung der Gesellschafter mittels eingeschriebener Briefe. Sie ist mit einer \emph{Frist
von mindestens einer Woche} zu bewirken. Wer eine offizielle Sitzung plant, muss solche
Fristen unbedingt einkalkulieren.
\end{beispiel}

\section{Die Sitzung vorbereiten}

Bevor die Einladung verfasst und versendet wird, sind mehrere \emph{Vorbereitungsaspekte} zu
bearbeiten und zu klären. Zu jedem Aspekt gehören Leitfragen, die vorab beantwortet werden
müssen.

\begin{center}
\begin{tabular}{p{4.2cm} p{9.0cm}}
\toprule
\textbf{Vorbereitungsaspekt} & \textbf{Mögliche Fragestellungen}\\
\midrule
Datum und Uhrzeit & Gibt es andere Sitzungen zum geplanten Termin? Sind Teilnehmer eventuell verhindert?\\
Raum reservieren und vorbereiten & Welcher Raum ist geeignet und verfügbar? Werden Medien/technische Hilfsmittel benötigt? Wurde der Raum reserviert?\\
Teilnehmer & Wer gehört zum Teilnehmerkreis? Welche Unterlagen benötigen die Teilnehmer zur Vorbereitung?\\
Zeitbedarf & Wie viel Zeit wird voraussichtlich für die Durchführung benötigt?\\
Tagesordnungspunkte (TOP) & Anmeldeschluss? Ziel je TOP? Zeit je TOP? Wer leitet ein? Medien nötig? Sinnvolle Reihenfolge?\\
Unterlagen & Welche Schriftstücke werden benötigt oder erhalten die Teilnehmer vorab?\\
Protokoll & Wer schreibt das Protokoll? Welche Protokollart wird gewünscht?\\
\bottomrule
\end{tabular}
\end{center}

\medskip
Neben Raum und Technik gehört auch die \emph{Verpflegung} (Getränke, bei längeren Sitzungen
Snacks) zu einer guten Vorbereitung --- sie hält die Konzentration aufrecht und signalisiert
Wertschätzung gegenüber den Teilnehmern.

\subsection{Die Tagesordnung und den Sitzungsablauf gestalten}

Sind die Fragen zu den Tagesordnungspunkten geklärt, wird der \emph{Sitzungsablauf}
strukturiert. Eine gewisse Wiederholung ist sinnvoll: Immer wiederkehrende Punkte (etwa
„Genehmigung des Protokolls“ oder „Verschiedenes“) können einen festen Platz im Ablauf
erhalten.

\begin{beispiel}
\textbf{Regelmäßiger Ablauf einer Abteilungsleitersitzung} (Uhrzeiten ca.):\\[0.3em]
09:00 Protokoll der letzten Sitzung \textbar{} 09:05 Aktuelle Berichterstattung aus den
Abteilungen \textbar{} 09:20 Schwerpunktthema \textbar{} 10:10 Pause \textbar{} 10:20 zwei
bis drei kürzere Themen \textbar{} 10:50 Ausblick, weitere Schritte, Verabredungen \textbar{}
11:00 Ende der Sitzung.
\end{beispiel}

Jeder Tagesordnungspunkt sollte ein klares \emph{Ziel} haben. Mögliche Ziele sind
Information, Ideensammlung, Erfahrungsaustausch, Aussprache, Entscheidungsfindung oder
Beschlussfassung. Vage Formulierungen sind gefährlich:

\begin{beispiel}
Ein TOP ist nur mit dem Wort „Sortiment“ bezeichnet. Damit bleibt unklar, was besprochen
werden soll: eine Erweiterung? Das Herausnehmen von Produkten? Eine Neuzusammenstellung? Die
Teilnehmer können sich nicht gezielt vorbereiten und gehen ohne durchdachte Vorschläge in die
Sitzung. Besser: „Aussprache und Entscheidung über die Sortimentserweiterung der Warengruppe
Handelswaren Bürotechnik (Anlagen 1 und 2)“.
\end{beispiel}

\subsection{Die Einladung erstellen}

Die Gestaltung der Einladung wird erst dann vorgenommen, wenn die Vorbereitungen weitgehend
abgeschlossen sind. \emph{Je präziser} die Inhalte in der Einladung formuliert werden, desto
besser können sich die Teilnehmer vorbereiten und Mitverantwortung für eine effektive
Durchführung übernehmen. Werden genaue Informationen erst in der Sitzung gegeben, beginnt das
Nachdenken erst dann --- und das gute Gelingen wird erschwert.

\begin{merksatz}[title={\bfseries Pflichtangaben einer detaillierten Einladung}]
\begin{itemize}[itemsep=1pt]
  \item Art der Sitzung/Besprechung
  \item Datum, Uhrzeit und Ort
  \item voraussichtlicher Zeitbedarf
  \item Tagesordnungspunkte mit Ziel (z.\,B.\ Information, Ideensammlung, Aussprache, Entscheidung, Beschlussfassung)
  \item Angabe des Berichterstatters/Verantwortlichen für einzelne TOP
  \item bei Bedarf Hinweis auf Anlagen zu einzelnen TOP
  \item Hinweis, ob Anmeldung erforderlich ist (für außerbetriebliche Teilnehmer)
  \item Name des Protokollführers (bei innerbetrieblichen Sitzungen)
\end{itemize}
\end{merksatz}

\subsection{Der zeitliche Ablauf der Vorbereitung (Countdown)}

Eine Checkliste verhindert, dass etwas vergessen wird. Der folgende Countdown zeigt einen
typischen Ablauf für eine größere Besprechung:

\begin{center}
\begin{tabular}{p{3.2cm} p{10.0cm}}
\toprule
\textbf{Zeitpunkt} & \textbf{Aufgabe}\\
\midrule
ca.\ 2 Wochen vorher & Ziel/Anlass klären, Teilnehmerkreis festlegen, Termin abstimmen\\
ca.\ 10 Tage vorher & Raum reservieren, Technik und Verpflegung klären\\
ca.\ 1 Woche vorher & TOP sammeln, Reihenfolge und Zeiten festlegen, Einladung versenden (Fristen beachten)\\
2--3 Tage vorher & Unterlagen/Anlagen zusammenstellen und versenden\\
am Sitzungstag & Raum und Technik prüfen, Moderationsmaterial bereitlegen, Getränke bereitstellen\\
\bottomrule
\end{tabular}
\end{center}

\begin{eselsbruecke}
\textbf{„Ziel --- Raum --- Runde --- Reihenfolge --- Einladung --- Unterlagen --- Check“.}
Diese Reihenfolge fasst den Vorbereitungs-Countdown zusammen: erst das \emph{Ziel}, dann
\emph{Raum} und Technik, der Teilnehmer\emph{kreis} (Runde), die \emph{Reihenfolge} der TOP,
die \emph{Einladung}, die \emph{Unterlagen} und am Ende der \emph{Check} vor Ort.
\end{eselsbruecke}

\section{Die Besprechung durchführen}

Eine durchdachte Vorbereitung ist ein wesentlicher Baustein für eine erfolgreiche Sitzung.
Dennoch kann es passieren, dass die Sitzung auch bei guter Planung nicht effektiv verläuft.
Denn für das gute Gelingen ist nicht nur die Leitung verantwortlich --- auch die
\emph{Teilnehmer} tragen Verantwortung.

\subsection{Rollen in der Besprechung}

\begin{itemize}
  \item \textbf{Leitung} --- eröffnet und schließt die Sitzung, achtet auf Ziel, Zeit und Beschlussfähigkeit, sorgt für Ergebnisse.
  \item \textbf{Moderation} --- steuert die Diskussion, erteilt das Wort, fasst zusammen, hält die Gruppe beim Thema und am Zeitplan (oft in Personalunion mit der Leitung).
  \item \textbf{Protokollführung} --- hält Verlauf, Ergebnisse und Beschlüsse fest; wird bereits in der Vorbereitung benannt.
\end{itemize}

\subsection{Schwierige Teilnehmertypen}

Verschiedene Persönlichkeiten werden in einer Sitzung sichtbar und können den Ablauf
ungünstig beeinflussen. Das Lehrbuch nennt drei typische Beispiele:

\begin{center}
\begin{tabular}{p{3.6cm} p{4.6cm} p{4.6cm}}
\toprule
\textbf{Typ} & \textbf{Verhalten} & \textbf{Reaktion der Moderation}\\
\midrule
Dauerredner & hat zu allem etwas zu sagen & freundlich begrenzen, auf Zeit und Thema verweisen\\
Desinteressierter & schaut immer wieder auf die Uhr & gezielt einbinden, nach seiner Meinung fragen\\
Abteilungswitzbold & macht auf Kosten anderer Scherze, zieht Beiträge ins Lächerliche & sachlich zum Thema zurückführen, Beiträge wertschätzen\\
\bottomrule
\end{tabular}
\end{center}

\subsection{Grundsätze und Verhaltensregeln}

Für die Durchführung von Sitzungen und Besprechungen helfen einige Grundsätze:

\begin{itemize}
  \item Die Sitzung beginnt \textbf{pünktlich}, auch wenn noch nicht alle Teilnehmer eingetroffen sind. Der vorgesehene Zeitrahmen wird eingehalten.
  \item Die \textbf{Tagesordnung} dient als Leitfaden; die Themen werden zielorientiert abgearbeitet. Die Teilnehmer schweifen nicht ab und führen keine Binnengespräche.
  \item Die Besprechung der einzelnen TOP folgt einem \textbf{strukturierten Ablauf}, z.\,B.\ Einführungs- und Informationsphase, Diskussionsphase, Entscheidungsphase --- oder kurz: \emph{fragen -- zuhören -- klären -- festhalten}.
  \item Die \textbf{Kommunikation} ist offen und respektvoll. Man fällt dem Redner nicht ins Wort, bemüht sich um Konsens, sieht sich als gleichwertige Gesprächsteilnehmer und spricht unterschiedliche Vorstellungen offen an.
  \item Jeder bemüht sich um \textbf{kurze Beiträge}.
\end{itemize}

\section{Elektronische und hybride Besprechungen}

Besprechungen beanspruchen viel Arbeitszeit. Deshalb planen immer mehr Unternehmen
alternative Formen, die nicht die Anwesenheit aller Betroffenen an einem Ort erfordern. Dazu
zählen vor allem \emph{Video-} und \emph{Telefonkonferenzen}.

\subsection{Video- und Telefonkonferenz}

Eine \textbf{Videokonferenz} ermöglicht, dass sich Gesprächspartner an weit
auseinanderliegenden Orten austauschen. Dafür werden ein spezielles Videokonferenzsystem und
ein schneller Internetanschluss benötigt. Sie erlaubt auch die gemeinsame Arbeit an einer
Datei bzw.\ den Austausch von Dateien, sodass man fast das Gefühl hat, nebeneinander zu
sitzen. So lassen sich Kosten (Flug-, Hotelkosten, Verpflegung) und Zeit einsparen; zunächst
fallen jedoch Anschaffungskosten für die Technik an.

Bei einer \textbf{Telefonkonferenz} wird ein Telefonat mit \emph{mehr als drei} Teilnehmern
geführt. Sind es genau drei, spricht man von einer \emph{Dreierkonferenz}. Es wird eine
spezielle Software benötigt; manche Anbieter stellen virtuelle Konferenzräume bereit, die nur
über eine PIN „betreten“ werden können und Platz für mehrere hundert Teilnehmer bieten.

\subsection{Moderne Praxis: Tools, Hybridmeetings und Etikette}

In der Praxis kommen heute Plattformen wie Microsoft Teams, Zoom oder Webex zum Einsatz; zur
Terminfindung dienen Online-Tools wie Doodle oder der Google Kalender. Häufig sind Meetings
\emph{hybrid}: Ein Teil sitzt im Konferenzraum, ein anderer ist online zugeschaltet. Damit
das funktioniert, gelten besondere Regeln.

\begin{merksatz}[title={\bfseries Etikette für Video- und Hybridkonferenzen}]
\begin{itemize}[itemsep=1pt]
  \item Pünktlich erscheinen und die Technik (Kamera, Mikrofon, Ton) vorher testen.
  \item Mikrofon stummschalten, solange man nicht spricht; Kamera möglichst einschalten.
  \item Ruhiger, aufgeräumter Hintergrund; Störquellen ausschalten.
  \item Eine klare Agenda und Moderation; nur eine Person spricht zur Zeit.
  \item Bei Hybridmeetings die Remote-Teilnehmer aktiv und gleichberechtigt einbinden.
\end{itemize}
\end{merksatz}

\section{Zusammenfassung}

\begin{enumerate}
  \item Gute \textbf{Vorbereitung} entscheidet über den Erfolg; mangelnde Planung führt zu Zeitverlust und vergessenen Beschlüssen.
  \item Die Vorbereitung beginnt \textbf{rechtzeitig} und beachtet \textbf{Fristen} (z.\,B.\ § 51 GmbHG: mind.\ 1 Woche).
  \item Zu klären sind \textbf{Vorbereitungsaspekte}: Datum/Uhrzeit, Raum, Teilnehmer, Zeitbedarf, TOP, Unterlagen, Protokoll.
  \item Die \textbf{Tagesordnung} braucht klare Ziele je TOP; vage Formulierungen verhindern eine gezielte Vorbereitung.
  \item Die \textbf{Einladung} enthält alle Pflichtangaben und wird erst nach Abschluss der Vorbereitung erstellt.
  \item Die \textbf{Durchführung} braucht Rollen (Leitung, Moderation, Protokoll), Pünktlichkeit, Zeit-/Themensteuerung und Diskussionsregeln.
  \item \textbf{Elektronische und hybride Besprechungen} sparen Zeit und Kosten, erfordern aber Technik, klare Moderation und Etikette.
\end{enumerate}

\newpage

\section*{Glossar}
\addcontentsline{toc}{section}{Glossar}

\glossareintrag{Sitzung/Besprechung}{Zusammenkunft mehrerer Personen zur Information, Aussprache oder Entscheidung über betriebliche Themen.}

\glossareintrag{Vorbereitungsaspekte}{Die vor einer Sitzung zu klärenden Punkte: Datum/Uhrzeit, Raum, Teilnehmer, Zeitbedarf, TOP, Unterlagen, Protokoll.}

\glossareintrag{Tagesordnung}{Geordnete Liste der in einer Sitzung zu behandelnden Themen (Tagesordnungspunkte, TOP), jeweils mit Ziel.}

\glossareintrag{Tagesordnungspunkt (TOP)}{Einzelner Punkt der Tagesordnung mit klarem Ziel, Verantwortlichem und eingeplanter Zeit.}

\glossareintrag{Einladung}{Schriftliche Aufforderung zur Teilnahme, die alle Pflichtangaben (Art, Datum/Ort, Zeitbedarf, TOP, Verantwortliche, Anlagen, Protokollführer) enthält.}

\glossareintrag{Einberufungsfrist}{Gesetzlich oder organisatorisch vorgegebener Mindestabstand zwischen Einladung und Sitzung (z.\,B.\ § 51 GmbHG: mind.\ eine Woche).}

\glossareintrag{Sitzungsablauf}{Strukturierte zeitliche Abfolge der TOP, oft mit festen, wiederkehrenden Punkten.}

\glossareintrag{Leitung}{Person, die die Sitzung eröffnet, leitet und schließt sowie auf Ziel, Zeit und Ergebnisse achtet.}

\glossareintrag{Moderation}{Steuerung der Diskussion: Wort erteilen, zusammenfassen, beim Thema und Zeitplan halten.}

\glossareintrag{Protokollführung}{Person, die Verlauf, Ergebnisse und Beschlüsse einer Sitzung schriftlich festhält.}

\glossareintrag{Videokonferenz}{Austausch von Gesprächspartnern an verschiedenen Orten per Videokonferenzsystem und schnellem Internet; ermöglicht gemeinsames Arbeiten an Dateien.}

\glossareintrag{Telefonkonferenz}{Telefonat mit mehr als drei Teilnehmern; spezielle Software, teils virtuelle Konferenzräume mit PIN.}

\glossareintrag{Dreierkonferenz}{Telefonkonferenz mit genau drei Teilnehmern.}

\glossareintrag{Hybridmeeting}{Besprechung, bei der ein Teil der Teilnehmer vor Ort und ein anderer Teil online zugeschaltet ist.}

\glossareintrag{Online-Tool zur Terminfindung}{Kostenloses Hilfsmittel (z.\,B.\ Doodle, Google Kalender) zur Abstimmung passender Termine.}

\end{document}
